\section{Momentum strategy}
\label{sec:momentum}

Momentum buys recent winners and shorts losers \citep{jegadeeshtitman1993}.
Our selected configuration, \#1292, uses six formation months and a one-month
skip. For prepared month-end close $P_{i,t}$, its signal is
\begin{equation}
 M_{i,t}=\frac{P_{i,t-1}}{P_{i,t-7}}-1.
 \label{eq:momentum}
\end{equation}

We rank eligible stocks in descending order for longs and ascending order for
shorts, breaking ties by \texttt{Asset\_ID}. Eligibility requires finite signals,
historical membership, sector records and valid price/event valuations.
Initially we take ten stocks on each side, with no overlap. Later decisions
retain incumbents within their side's top fifteen ranks ($10\times1.5$) and
fill vacancies in rank order.

\subsection{Sizing and trading}

With gross exposure $G=1.5$ and long share $a=0.7$, target weights are
\begin{equation}
 w_{i,t}^{\mathrm{raw}}=
 \begin{cases}
   Ga/10=+0.105, & i\in L_t,\\
   -G(1-a)/10=-0.045, & i\in S_t,\\
   0, & \text{otherwise}.
 \end{cases}
 \label{eq:weights}
\end{equation}

The book is 105\% long and 45\% short: 150\% gross exposure and 60\% net long exposure.
Positions are equal within each side, with no sector neutrality.
We impose no fixed single-position cap.

Targets convert to whole units on each domain's declared price basis, using
the sizing price and NAV. Same-side resizing is suppressed when absolute
weight drift is at most 0.05 of NAV. Entries, exits and sign changes still
trade; feasibility repairs can override suppression. This five-percentage-point
threshold exceeds a short position's 4.5\% target, allowing substantial weight
drift.

The historical engine decides, sizes and trades at the labelled month-end
close, then values the next interval. The signal is lagged, but sizing and
spread inputs use that close. The competition replay instead executes
month-end targets at the next scheduled trading-day open.

The rank buffer and resize threshold limit small adjustments, not losses.
Concentrated longs, sector imbalance, short squeezes and leverage remain
material risks.
